\section{Related Work}\label{sec:related}
Choosing among several deployed AI coding agents to maximize task success net of cost relates to five lines of work: 1) cost-aware LLM routing, 2) AI coding agent evaluation, 3) empirical studies of agentic pull requests, 4) off-policy evaluation, and 5) pitfalls of mining software repositories. We review each and contrast it with our setting, in which the candidates are whole commercial agents and the only signal is observational GitHub outcomes.

\subsection{Cost-Aware LLM Routing}

LLM routing addresses the cost-quality trade-off by using a learned router to decide \emph{a priori} which model should handle a given query. \textbf{Zooter}~\cite{lu2023routing} introduced reward-guided routing that distills rewards from training queries into a routing function, achieving top performance on 44\% of tasks across 26 benchmark subsets. \textbf{Hybrid-LLM}~\cite{ding2024hybrid} proposed quality-aware routing that trades quality for cost based on predicted query difficulty, achieving up to 40\% fewer calls to the large model with no quality degradation. \textbf{RouterBench}~\cite{hu2024routerbench} established a standardized framework of over 405k inference outcomes for comparing routing strategies. \textbf{RouteLLM}~\cite{ong2024routellm} learns routers from human preference data, reporting cost reductions over 85\% on MT-Bench and 35\% on GSM8K while retaining 95\% of GPT-4's performance. Most recently, \textbf{GraphRouter}~\cite{GraphRouter} models tasks, queries, and LLMs as a heterogeneous graph and predicts performance-cost over edges, improving at least 12.3\% over prior routers and generalizing to unseen LLMs. A related cost-quality line cascades models rather than routing once: \textbf{FrugalGPT}~\cite{chen2024frugalgpt} escalates through an LLM cascade and reports up to 98\% cost reduction at matched accuracy.

These methods share two structural assumptions that do not hold in our setting. First, they route over a \emph{shared input}: the same query can in principle be sent to any candidate, and preference-labeled or overlapping responses provide a near-counterfactual training signal. Second, they choose among \emph{models} and are evaluated on short-form benchmarks with graded correctness. When the candidates are whole deployed coding agents and the only data is observational, each task is attempted by exactly one agent with essentially no overlap, the success label (merge) is a workflow-confounded proxy rather than a graded answer, and, as we show, the per-task success ranking is not even identifiable (Section~\ref{sec:findings}).

\subsection{AI Coding Agents and Their Evaluation}

A parallel line evaluates coding agents directly. \textbf{SWE-bench}~\cite{jimenez2024swebench} measures whether a system resolves real GitHub issues and has become the standard counterfactual benchmark, with leaderboard infrastructure such as the \textbf{Holistic Agent Leaderboard}~\cite{holistic2025leaderboard} running many systems on the same curated tasks. Cost is now a first-class concern: \textbf{EET}~\cite{eet2026} cuts the spend of a \emph{single} software-engineering agent by terminating unpromising trajectories early, and \textbf{task-level evaluations}~\cite{tasklevel2026eval} compare agents head-to-head on shared open-source tasks. The \textbf{SE~3.0} vision~\cite{li2025se3} frames agents as first-class teammates, and field experiments~\cite{cui2024genai} estimate the productivity effect of generative AI on developers with a randomized design.

All of these either optimize cost \emph{within} a single agent or rank agents on \emph{shared} tasks that supply counterfactual overlap. None studies the operational decision of which deployed agent to assign to an incoming task when the only evidence is each agent's past, non-overlapping behavior in the wild. Our work targets exactly that decision.

\subsection{Empirical Studies on Agentic Pull Requests}

The \textbf{AIDev} dataset~\cite{li2026aidev}, 932,791 agent-authored pull requests across five commercial agents and 116k repositories, has driven a wave of descriptive studies: the security of agent-authored PRs~\cite{security2026teammates}, whether agents contribute tests~\cite{haque2026tests}, why agentic fixes are rejected~\cite{rejection2026fixes}, and task-level agent performance in open-source projects~\cite{tasklevel2026eval}.

This body of work characterizes \emph{what} agents do and \emph{how well} they do it, but treats agent identity as a \emph{descriptive} variable. We instead treat the agent as the quantity to be \emph{chosen}, turning the dataset from a description into a decision problem. Doing so immediately exposes a confound the descriptive studies do not have to confront: the agent that looks best is largely the one deployed in the easiest repositories under a self-merge workflow (Section~\ref{sec:findings}). This step, from describing agents to routing among them under cost, is the main way our scope extends the AIDev studies we build on.

\subsection{Off-Policy Evaluation}

Because each task is attempted by a single agent, valuing a routing policy is an off-policy evaluation (OPE) problem. The standard estimators are the direct method (fit an outcome model and average its predictions under the new policy), inverse propensity weighting (IPW), and the doubly robust combination~\cite{dudik2011doubly}, developed for contextual-bandit policy evaluation. We adopt a subgroup direct-method estimator and cross-check it with self-normalized IPW, a doubly robust estimator over a covariate propensity model, and a two-level bootstrap.

OPE is normally applied where the logging policy is at least partially stochastic and overlap is plausible. Our logging policy (how developers picked agents) is neither randomized nor overlapping, so even doubly robust estimates rest on a no-unobserved-confounding assumption that observational data cannot test; we therefore report the \emph{disagreement} among estimators as a first-class result rather than a single point estimate.

\subsection{Mining-Repository Pitfalls and Outcome Prediction}

A methodologically adjacent line predicts per-commit or per-PR outcomes: \textbf{just-in-time defect prediction}~\cite{ni2024jit} forecasts defects without modeling which author produced the change, and pull-request acceptance prediction is a related task; class imbalance in such labels is itself a studied problem~\cite{hegarcia2009imbalanced}. A broader literature documents the hazards of treating repository signals at face value: \textbf{Kalliamvakou et al.}~\cite{kalliamvakou2016promises} catalog the promises and perils of mining GitHub, and \textbf{Menzies and Shepperd}~\cite{menzies2019badsmells} catalog analogous ``bad smells'' in software-analytics studies.

Our contribution to this thread is to add agent identity as a feature and show that, once it is present, merge behaves as a noisy and confounded label dominated by per-agent base rates rather than task content. The resulting selection bias is a concrete, quantified instance of the sampling-assumption hazards these works warn about.

\subsection{Positioning of Cost-Aware AgentRouter}

Our work differs from all of the above. Unlike LLM routing and cascading, which choose among \emph{models} over a shared input with near-counterfactual or benchmark labels, we choose among whole deployed \emph{agents} from purely observational outcomes; the consequence is that the success-axis ranking is not identifiable, so we relocate the contribution to the cost axis. Unlike the AIDev empirical studies, which describe agents, we treat the agent as a decision variable. Unlike standard off-policy evaluation, which assumes some overlap, we operate with one agent per task and report estimator disagreement as evidence rather than hiding it behind a point estimate.

To our knowledge, this is the first work to frame the selection among deployed commercial AI coding agents as a cost-aware routing problem learned and evaluated on purely observational GitHub data, and to show that on such data the per-task success ranking is not identifiable while a robust, cost-saving routing policy nonetheless exists.
